\section*{Exercises}
\problem{
{\it Wilson loops\aidx{Wilson loop}}
% See Appendix A of Tomboulis & Yaffe  Comm Math Phys 100, 313, 1985 Finite Temp SU2 gauage theory

Consider a reflection $\Theta$ that maps two equal partitions of a lattice geometry, $\Lambda_{\pm}$, into each other: $\Theta(\Lambda_\pm)=\Lambda_\mp$. A reflection positive measure guarantees that for any quantity $F$ that has support only on the degrees of freedom in $\Lambda_+$, $\langle F\Theta(F)\rangle\ge0$.
\begin{enumerate}
    \item[(a)] Show that for two observables $F$ and $G$ with support in $\Lambda_+$, reflection positivity\aidx{reflection positivity} implies that 
    $$
    |\langle F \Theta(G)\rangle|\leq \sqrt{\langle F\Theta(F)\rangle \langle G\Theta(G)\rangle}
    $$
\end{enumerate}

Now consider rectangular Wilson loops $W_{I,J}$ of size $I\times J$ in pure Yang-Mills theory. 
\begin{enumerate}
    \item[(b)] Use reflection positivity to show that  $W_{1,1} \leq {W_{1,2}}^{1/2}$.
    \item[(c)] Extend your argument to prove that $W_{1,1} \leq {W_{I,J}}^{1/IJ}$ for all $I,J\in \mathbb{Z}^+$.
    \item[(d)] Discuss how this provides an upper bound on the rate of growth of the static quark potential\aidx{static quark potential}
    $$V(R)=-\lim_{T\to\infty}W_{R,T}/T.$$
\end{enumerate}
}
{
First argue that $\langle \ldots,\Theta(\ldots)\rangle$ is an inner product on $\Lambda_+$ because of reflection positivity. Then Cauchy-Schwartz immediately implies (a). 
Then set F=product of 3 sides of plaquette\aidx{plaquette} and G=4th side. Then for reflection in plane containing G, $\Theta(G)=G^\dagger$ so $G\Theta(G)=1$ and $F\Theta(F)=W_{2,1}$ is 2x1 loop
 Now repeat this procedure for 2x1 to bound 4x1 etc. 
 For J not multiple of 2, use eg 2x1 loop reflected at 1.5a to give 
\begin{eqnarray}
    W_{12} \leq W_{13}^{1/2} W_{11}^{1/2} \\
    W_{11}\leq W_{12}^{1/2} \leq W_{13}^{1/4} W_{11}^{1/4}\\
    \to W_{11}\leq W_{13}^{1/3}
\end{eqnarray}
Repeat argument in the orthogonal direction of the loop to bound $W_{IJ}$
The results bound V to grow at most linearly with R
}


\problem{{\it Parallel transporters}
    
The solution $E\left(t, t_0\right)$ to the differential equation
$$
\left[\frac{d}{d t}+M(t)\right] E\left(t, t_0\right)=0
$$
with $M(t)^\dagger=-M(t)$ and
$$
\left[M\left(t_1\right), M\left(t_2\right)\right] \neq 0
$$
is given by
$$
E\left(t, t_0\right)=\mathcal{P} \exp \left[-\int_{t_0}^t d s M(s)\right].
$$
Prove the following statements
\begin{enumerate}
    \item 
    $\frac{d}{d t_0} E\left(t, t_0\right)-E\left(t, t_0\right) M\left(t_0\right)=0$
\item $E(t, u) E(u, v)=E(t, v)$
\item $E(t, u)=E^{-1}(u, t)$
\item  $E^{\dagger}(t, u)=E^{-1}(t, u)$
%\item Under a gauge transformation

\end{enumerate} 
    }
    {224}

\problem{{\it Haar measure\aidx{Haar measure}}

    The Haar measure on a group G with elements $U\in G$ parameterised as $U=U(\omega)$ for $\omega=(\omega_1,\ldots,\omega_n)$ is
    $$
    dU = h(\omega) d\omega_1\ldots d\omega_n
    $$
    with 
    $h(z)=c\sqrt{\det(g(\omega))}$
    where $c$ is a normalisation constant and 
    $$
    g(\omega)_{ij}={\rm tr}\left[\frac{\partial U}{\partial\omega_i}\frac{\partial U^\dagger}{\partial\omega_j}\right].
    $$
    \begin{enumerate}
        \item Show that for a well-behave function $f(U)$ and a Haar measure $dU$
        $$
        \int dU f(U) = \int dU f(U^\ast) = \int dU f(U^{-1})
        $$

        \item For $G=SU(2)$
        \begin{enumerate}
            \item Consider the parameterisation of $U\in SU(2)$ in terms of coordinates on the 3-sphere $S^3\subset R^4$
            $$
            U=x_0 + i x_a \sigma_a,\qquad x^2=\sum_{a=0}^3 x_a^2 =1,
            $$
            where $\sigma_a$ are Pauli matrices.  Construct the Haar measure in these coordinates

            \item Now instead use stereographically projected coordinates in ${\bf z}=(z_1,z_2,z_3) \in R^3$
            where 
            $$x_0=\frac{1-|{\bf z}|^2}{1+|{\bf z}|^2}, \qquad x_a=\frac{2 z_a}{1+|{\bf z}|^2}
            $$
            and construct the Haar measure.

            \item Construct the Haar measure for the parameterisation $U=\exp(i\sigma_a \theta_a)$.
            \end{enumerate}
        
    \end{enumerate}
}
{211}
    
\problem{{\it Gauge actions for SU(N)}
\begin{enumerate}
    \item Show that the Wilson gauge action\aidx{Wilson gauge action} 
    $$
    S_W=-\frac{\beta}{2N}\sum_{P_{\mu\nu}}\{ {\rm tr}[P_{\mu\nu}]+{\rm tr}[P^{-1}_{\mu\nu}]\}
    $$
    (where the sum is over all plaquettes) reduces to the continuum gauge action in the geometric $a\to0$ limit.
    
    \item Consider the Symanzik action
    $$
    S_I=-\frac{\hat \beta}{2N}\left[\sum_{P_{\mu\nu}}\{ {\rm tr}[P_{\mu\nu}]+{\rm tr}[P^{-1}_{\mu\nu}]\} + \gamma \sum_{R_{\mu\mu\nu}}\{ {\rm tr}[R_{\mu\mu\nu}]+{\rm tr}[R^{-1}_{\mu\mu\nu}]\}\right],
    $$
    where the sums are over all plaquettes and over all $2\times1$ ($\mu\ne\nu$) rectangles 
    $$
        R_{\mu\mu\nu}= U_\mu(n) U_\mu(n+\hat\mu) U_\nu(n+2\hat\mu)U^{-1}_\mu(n+\hat\mu+\hat\nu)U^{-1}_\mu(n+\hat\nu)U^{-1}_\nu(n).
    $$
    Show that the correct continuum limit\aidx{continuum limit} is obtained. Find $\gamma$ and $\hat\beta$ such that the leading discretisation artifacts\aidx{improvement, $\mathcal{O}(a)$}\aidx{improvement, $\mathcal{O}(a)$} are ${\cal O}(a^4)$ and find the relation between $\hat\beta$ and $\beta$ in the unimproved action.
    
    \item Consider the Manton action
    $$
    S_M=-\frac{\tilde\beta}{N}\sum_{P_{\mu\nu}} [d(P_{\mu\nu},1)]^2 ,
    $$
    where $d(P,Q)$ is the \textit{shortest} distance between two elements in the gauge group $G$  and show this has the correct continuum limit. 
    It can be shown that $d(P,Q)={\rm tr}[\chi^2]^{1/2}$ where $\chi$ is an element of the Lie algebra $\mathfrak{g}$ such that $PQ^{-1}=\exp(i\chi)$ with all eigenvalues $\lambda_i\in(-\pi,\pi)$ (since $G$ is compact, there are multiple algebra elements that generate a given group element by exponentiation, but the latter condition selects geodesic path with the shortest distance).
    
    Show that this action has the same continuum limit as the Wilson action for an appropriate coupling $\tilde\beta$.
    
    [{\sc Hint:} see Manton, Phys. Lett.B 96 (1980) 328 for context]
\end{enumerate}
}    
{112}
